\begin{table}[H]
\centering
\scriptsize
\caption{Findings register. Severities are PROPOSED and must be confirmed by the validator.}
\label{tab:T7_findings}
\begin{tabular}{>{\raggedright\arraybackslash}p{0.5cm}>{\raggedright\arraybackslash}p{5.2cm}>{\raggedright\arraybackslash}p{1.8cm}>{\raggedright\arraybackslash}p{2.7cm}>{\raggedright\arraybackslash}p{3.5cm}}
\toprule
ID & Finding & Severity & Evidence & Proposed remediation \\
\midrule
F1 & With a 250-day window the correctly specified GARCH(1,1)-t over-rejects: Kupiec 10.6 $\pm$ 1.0\,\% vs exact size 5.5\,\%, Z2 17.7 $\pm$ 1.2\,\% (T = 1000); exceedance rate 1.25\,\%. With 1000 days: Kupiec 5.8 $\pm$ 0.7\,\%. Estimation-window study: the rejection rates fall to about the reference size by 500 days (GARCH-t) and 750 days (FHS) (Kupiec, \% by window: 250: 11.8, 500: 4.5, 750: 3.6, 1000: 3.7), but the VaR stays slightly anti-conservative: exceedance rate 1.05--1.13\,\% against the true VaR's 1.01\,\% (GARCH-t by window: 250: 1.27, 500: 1.11, 750: 1.06, 1000: 1.05). Paired difference against the true VaR at 1000 days: +0.041 $\pm$ 0.007 pp. & PROPOSED: High & Tab.~\ref{tab:T3_power}, \ref{tab:T4_estimation_window}; Fig.~\ref{fig:fig3_estimation_window} & Estimation window of at least 500 days for GARCH-t and 750 days for FHS (where the rejection rates reach the reference sizes); account for the remaining small anti-conservatism (parameter uncertainty), e.g. by a VaR add-on or a longer window. \\
F2 & FHS with a 250-day window underestimates the tail (ES): Z2 rejects in 33.6 $\pm$ 1.5\,\% (oracle 4.6\,\%, T = 1000); with 750 days 6.3 $\pm$ 0.8\,\%. & PROPOSED: Medium & Tab.~\ref{tab:T3_power}, \ref{tab:T4_estimation_window} & Longer window for the residual distribution; monitor the ES backtest separately from VaR. \\
F3 & Static models miss volatility clustering: the static normal VaR on GARCH-t data is rejected by cc in 59.7 $\pm$ 1.6\,\% and by Z2 in 64.7 $\pm$ 1.5\,\% (T = 1000). & PROPOSED: High & Fig.~\ref{fig:fig1_exceedances}, \ref{fig:fig2_rejection_rates}; Tab.~\ref{tab:T3_power} & Use a conditional volatility model (GARCH-t / FHS with an adequate window). \\
F4 & The Kupiec test is mis-sized at T = 250: exact size 9.48\,\% at a nominal 5\,\% (x = 0 is rejected); the exact binomial test keeps 4.12\,\%. & PROPOSED: Medium & Tab.~\ref{tab:T2_size} & Report the exact binomial test (or the exact size) alongside Kupiec for one-year backtests. \\
F5 & The Christoffersen independence test is undersized (0.9--2.6\,\% on correct VaRs) and has little power on its own. & PROPOSED: Low & Results file of the power study (ind columns) & Do not rely on the independence test alone; use cc and the ES test. \\
F6 & $\sqrt{10}$ scaling of a fat-tailed one-day VaR overstates the 10-day VaR on iid t5 data (+8.6 $\pm$ 0.3\,\%). On GARCH-t6 data $\sqrt{10}$ is unbiased on average but $\sqrt{10}\times$ static normal has an RMSE of 29.3\,\%. & PROPOSED: Medium & Fig.~\ref{fig:fig5_horizon}; Tab.~\ref{tab:T5_horizon} & Simulate the 10-day distribution (GARCH Monte Carlo); use $\sqrt{10}$ only for a conditional one-day VaR when the 10-day sum is near-normal. \\
F7 & The stationary block bootstrap is biased on GARCH-t6 data (+3.7 $\pm$ 1.2\,\%, RMSE 37.2\,\%): it does not condition on the current volatility. & PROPOSED: Low & Tab.~\ref{tab:T5_horizon} & Do not use the unconditional bootstrap as the 10-day benchmark under volatility clustering. \\
F8 & Historical simulation with $k = \lceil n\alpha\rceil$ has a finite-sample exceedance probability of $(n-k+1)/(n+1)$ = 1.20\,\% at n = 250, $\alpha$ = 0.99 (simulated: 1.21\,\% on iid normal data). & PROPOSED: Low & Results of the power study (HS on iid normal data) & Document the quantile convention; consider a longer window or interpolation. \\
F9 & Survivorship (synthetic, hazard strength 1, CAGR $\ge$ 10\,\%): a survivors-only database shows 34.3\,\%, against 25.5\,\% for all firms had they survived (selection bias 8.8 $\pm$ 0.02 pp) and 17.8\,\% for the full cohort with exits as misses (attrition 7.8 $\pm$ 0.01 pp). The size depends on the assumed hazard; the $\pm$ se is Monte-Carlo noise only. & PROPOSED: Medium & Fig.~\ref{fig:fig7_survivorship}; Tab.~\ref{tab:T6_survivorship} & Use point-in-time data that include delisted firms, and state which base rate (counterfactual or full cohort) answers the question asked. \\
\bottomrule
\end{tabular}
\end{table}
